\section{Transparency: can I see why it did that?}
\label{sec:transparency}

A clinician cannot endorse a black box. The third question is whether the system's
reasoning is inspectable and whether a complete, tamper-evident record lets the
clinician reconstruct events after the fact.

[DRAFTING INSTRUCTIONS: in the full-framework, describe explainable output and the
hash-chained audit trail, grounded in the FDA action plan
\texttt{fda2021samd} and the author's records and server lineage in
\texttt{review/references/author\_works.bib}
(\texttt{Kawchak2026NationalMCPServersPhysical},
\texttt{Kawchak2026PhysicalAIOncologyClinical}). Reproduce Mermaid
\texttt{adoption/mermaid/diagrams/08-audit-trail-sequence.md} and
\texttt{18-tumor-board-defensibility-mindmap.md} as colored cfwfig TikZ. Carry the
KEY TERM ``audit trail''. Build a body-width table of what the record contains and
who can read it.]

\begin{cfwfig}
\node[cftwo] (a) at (0,0) {Action}; \node[cfact] (b) at (3.6,0) {Hash-chained record};
\node[cfhope] (c) at (7.6,0) {Auditor reconstructs};
\draw[cfedge] (a)--(b); \draw[cfedge] (b)--(c);
\draw[cfedged] (b) to[out=120,in=60] (a.north);
\end{cfwfig}
\vspace{-0.2cm}
\cfwcaption{Figure 5. The audit trail (placeholder; expanded from Mermaid 08 in
the full-framework).}

Transparency is operational, not rhetorical: the test is whether an auditor can
reconstruct exactly what happened, and why, from the record alone.
